\section*{Conclusion} \label{sec:conclusion}
This report has defined and scoped a vision-based docking system for an inspection-class ROV approaching a free-floating underwater station, motivated by the need to reduce reliance on risky, vessel-dependent launch and recovery.

Through a literature review, it has surveyed acoustic, electromagnetic, mechanical and vision-based docking methods, and identified key gaps: most work assumes seabed-mounted or otherwise static docks, focuses on AUVs rather than tethered ROVs, and is validated mainly in clear, controlled conditions rather than turbid, dynamic environments. In response, the project has been framed around a turbidity-adaptive, marker-based pose estimation pipeline for a BlueROV-style platform, with a proposed architecture, aiming to enhance the feasibility for low cost autonomous underwater opertions.

Within Thesis A, foundational preparations have been completed: the software environment and ROS 2 Humble stack have been configured, camera calibration performed, and a working ArUco-based pose estimation node and initial dock CAD concepts implemented. Next Thesis B and C work will extend this prototype to underwater imagery and turbidity-aware pre-processing, integrate it with the ROV's navigation and control, and experimentally evaluate autonomous docking performance under representative conditions.